% sections/01_introduction.tex

\section{Introduction}
\label{sec:introduction}

\subsection{Chemometrics Context}

Near-infrared (NIR) and Fourier-transform infrared (FTIR) spectroscopy are
workhorses of industrial and pharmaceutical quality control, providing fast,
non-destructive, reagent-free characterization of chemical composition and
physical properties~\citep{pasquini2003nir}.
Both modalities produce high-dimensional spectra ($p = 200$--$3000$ variables)
from small sample sets ($n = 10$--$200$ in most practical calibration
exercises)~\citep{martens1989multivariate,rinnan2009review}.

A structural feature of spectroscopic datasets that distinguishes them from
general ML benchmarks is \emph{replicate nesting}: many industrial and
laboratory protocols collect multiple spectral scans from the same physical
sample or production batch.
When a dataset contains, say, 5 replicate scans per sample and 20 samples,
naive $k$-fold cross-validation assigns replicates of the same sample to both
the training and test fold---inflating reported accuracy, sometimes to
perfect scores, while the model has learned nothing generalizable beyond the
instrument noise it memorized.
The correct approach, GroupKFold or LeaveOneGroupOut (LOGO) keyed on sample
identity, keeps all replicates of a sample on the same side of the
train/test boundary~\citep{kaufman2012leakage}.
In practice, this failure mode is pervasive and hard to detect automatically
without explicit replicate-structure parsing.

\subsection{Failure Modes of Ungrounded LLM Analysis}

Large language models (LLMs) have been proposed as general-purpose scientific
analysis agents~\citep{yao2023react,schick2023toolformer}.
Applied naively to spectroscopic data, they exhibit a characteristic cluster
of failure modes:

\begin{itemize}
  \item \textbf{Leakage blindness.} Without explicit group-structure checking,
    an LLM agent will accept and report a leakage-inflated cross-validation
    score as valid. The inflated metric then propagates into interpretation and
    recommendations.
  \item \textbf{Metric hallucination.} LLMs tend to produce plausible-sounding
    numeric metrics from prior knowledge rather than from computed outputs.
    A model that has never seen the dataset can still generate a ``typical''
    R\textsuperscript{2} for a wear-layer regression problem.
  \item \textbf{Causal overclaiming.} Feature importance in a regression or
    classification model does not imply causal mechanism.
    LLMs routinely conflate high SHAP values or regression coefficients with
    causal attribution unless explicitly constrained.
  \item \textbf{Silent fallback acceptance.} When a model fails to converge,
    an autonomous agent may silently substitute a fallback model and continue
    reporting results without flagging that the scientific plan has changed.
\end{itemize}

These failure modes are not specific to one LLM family; they arise structurally
from the mismatch between language-model pretraining (text prediction) and
the epistemological requirements of scientific reproducibility~\citep{kapoor2023leakage}.

\subsection{Thesis: Reliability from Structure, Not the Model}

The central thesis of this work is that trustworthy spectroscopic ML analysis
does not require a more capable or larger language model---it requires
architectural constraints that make errors detectable and gateable regardless
of which model drives the conversation.

We pursue three mechanisms:
(1)~\emph{typed contracts} that make every tool input and output a
serializable, schema-validated dataclass, preventing silent type coercions and
making all information explicit;
(2)~\emph{deterministic, MCP-free core logic} that produces identical outputs
for identical inputs at any seed, enabling byte-reproducible audits; and
(3)~\emph{HITL gates} that route non-routine decisions (plan approval,
fallback acceptance, final model selection under warnings) through a human
checkpoint rather than resolving them autonomously.

\subsection{Contributions}

This paper makes the following contributions:

\begin{enumerate}
  \item \textbf{MCP chemometrics server.} An open-source MCP server exposing
    8 core spectroscopic analysis tools and 3 persistent method-memory tools,
    designed for use with any MCP-compatible LLM client.

  \item \textbf{Typed contract layer.} A library of 18 domain contracts and
    12 request/response wrappers (30 frozen dataclasses) over a shared
    SerializableContract base, enforcing schema correctness across the tool boundary.

  \item \textbf{Declarative modality registry.} A frozen \texttt{ModalityProfile}
    registry (FTIR, NIR, Raman) that centralizes axis-unit enforcement,
    value-domain clipping, preprocessing scope, and modality inference from
    axis range---eliminating scattered \texttt{if modality == \textquotedbl{}FTIR\textquotedbl{}} branches
    from all downstream modules.

  \item \textbf{Leakage detection as a gating tool.} A \texttt{validate\_results}
    tool that automatically detects replicate leakage (via prediction-consistency
    within groups) and group leakage risk (when candidate group columns exist
    but grouped CV was not used), and gates downstream model selection on the
    resulting warnings.

  \item \textbf{Deterministic evaluation framework.} A Layer-1 harness
    (seed = 42, no LLM, direct tool import) that measures leakage detection,
    hallucination guards, fallback correctness, plan quality, and HITL
    compliance with byte-reproducible outputs across four real spectroscopic
    scenarios.
\end{enumerate}
